\section{The BFS Region-Growing Algorithm}
\label{sec:algorithm}

\subsection{Notation and Subgraph Setup}

Let $G_v = (V_v, E_v, w_v)$ denote the subgraph at bisection tree node $v$,
where $V_v$ is the set of census tracts assigned to the subtree rooted at $v$,
$E_v \subseteq E$ is the induced edge set, and $w_v$ inherits geographic
boundary-length weights from the parent graph (B.2 default).
Let $n = |V_v|$ be the number of tracts.
Let $p_i$ denote the population of tract $i \in V_v$, and let
$P_v = \sum_{i \in V_v} p_i$ be the total subgraph population.
For the bisection case, $k = 2$; the algorithm generalises directly to
$k > 2$ by growing $k$ regions simultaneously.

A \emph{bisection} at node $v$ is a partition $(L, R)$ of $V_v$ satisfying
the population-balance constraint:
\[
  \left|\sum_{i \in L} p_i - \frac{P_v}{2}\right|
  \;\le\;
  \delta \cdot \frac{P_v}{2},
\]
where $\delta = 0.005$ is the $\pm 0.5\%$ federal tolerance.

\subsection{Seed Derivation}

All randomness in \bfsg{} enters through seed selection.
The seed is derived deterministically from the global base seed:

\begin{definition}[BFS Growth Seed]
\label{def:bfs-seed}
Let $\mathtt{base\_seed}$ be the unsigned 64-bit integer global seed.
The \emph{BFS Growth seed} is
\[
  s_{\mathrm{bfs}} = \mathrm{first8}\bigl(
    \mathrm{SHA\text{-}256}(
      \texttt{"BFS\_SEED\_"} \;\|\; \mathrm{le8}(\mathtt{base\_seed})
    )
  \bigr),
\]
where $\mathrm{le8}(x)$ encodes $x$ as 8 little-endian bytes and
$\mathrm{first8}(h)$ reads the first 8 bytes of hash $h$ as a
little-endian unsigned 64-bit integer.
The derived PRNG is $\mathrm{SmallRng::seed\_from\_u64}(s_{\mathrm{bfs}})$.
\end{definition}

The domain-separation prefix \texttt{"BFS\_SEED\_"} embeds algorithm
identity in the seed derivation.
A test asserts that the prefix constant in source equals exactly
\texttt{"BFS\_SEED\_"} and that $s_{\mathrm{bfs}}(0)$ equals a
hard-coded expected value, preventing silent seed-compatibility drift.
Same $\mathtt{base\_seed} \Rightarrow$ identical seed placement
$\Rightarrow$ identical BFS assignment $\Rightarrow$ identical final plan.

\subsection{Seed Selection: k-Farthest via BFS}

\begin{definition}[k-Farthest Seeds]
\label{def:seeds}
Given sorted tract indices $\mathrm{sorted} = (i_0, i_1, \ldots, i_{n-1})$
and derived PRNG $\mathrm{rng}$:
\begin{enumerate}
\item \textbf{Seed~0.} Sample $\mathrm{seed}_0$ from $V_v$ by
  population-weighted sampling using $\mathrm{rng}$.
  Tracts with larger population have proportionally higher probability
  of being selected as seed~0.

\item \textbf{Seeds $1, \ldots, k-1$.} For $j = 1, \ldots, k-1$:
  compute BFS distances $\mathrm{dist}_i = \min_{l < j} d_{\mathrm{BFS}}(i, \mathrm{seed}_l)$
  for all unselected tracts $i$.
  Set $\mathrm{seed}_j = \arg\max_i \mathrm{dist}_i$.
\end{enumerate}
\end{definition}

The BFS-farthest selection ensures that seeds are maximally spread in
the graph topology.
For $k = 2$, seed~1 is the tract farthest from seed~0 in BFS hops,
which is the tract at the opposite end of the state's geographic extent.

\begin{remark}
Population-weighting for seed~0 reflects the observation that
high-population areas (cities, suburban cores) are more likely to
require splitting in the bisection tree.
Starting from a high-population region ensures that the first district's
growth immediately encounters the densest part of the graph, where
boundary decisions are most consequential.
If seed~0 were placed in a low-population rural area, the growing district
would accumulate many low-population tracts before reaching the urban
core, at which point the opposing district would already control most of
the dense tracts.
\end{remark}

\subsection{BFS Growth Phase}

\begin{definition}[Priority Queue Growth]
\label{def:growth}
Given seeds $(\mathrm{seed}_0, \ldots, \mathrm{seed}_{k-1})$
and current district populations $(P_0, \ldots, P_{k-1})$ initialised to
$(p_{\mathrm{seed}_0}, \ldots, p_{\mathrm{seed}_{k-1}})$:
\begin{enumerate}
\item Initialise assignment $A: V_v \to \{0, \ldots, k-1, \bot\}$ with
  $A(\mathrm{seed}_d) = d$ for $d = 0, \ldots, k-1$ and $A(i) = \bot$
  for all other tracts.

\item Initialise minimum-priority queue $\mathcal{Q}$ with entries
  $(P_d, \mathrm{seed}_d\text{'s neighbours}, d)$ for each seed $d$,
  where $P_d$ is the priority key.

\item While $\mathcal{Q}$ is non-empty:
  \begin{enumerate}
  \item Pop $(P_d, \mathrm{tract}\; i, d) = \mathcal{Q}.\mathrm{pop\_min}()$.
  \item If $A(i) \neq \bot$: skip (already assigned by another path).
  \item Set $A(i) = d$; update $P_d \mathrel{+}= p_i$.
  \item For each neighbour $j$ of $i$ with $A(j) = \bot$:
    push $(P_d, j, d)$ to $\mathcal{Q}$.
  \end{enumerate}
\end{enumerate}
\end{definition}

The priority $P_d$ at the time of each push reflects the district's current
population.
Because population is monotonically non-decreasing as tracts are assigned,
tracts pushed later for the same district have weakly higher priority keys,
so older entries near a smaller district always appear earlier in the
queue.
This ensures that the emptiest district expands first.

\subsection{Post-Hoc Rebalancing}

The BFS growth phase produces a contiguous assignment but does not
guarantee the $\pm 0.5\%$ population balance constraint.
Population imbalance arises when the final few tracts assigned to a
district push it over the ideal.
A post-hoc boundary-swap rebalance corrects this:

\begin{definition}[Rebalance]
\label{def:rebalance}
The rebalance phase performs up to \texttt{balance\_post\_iters} $= 200$
boundary-swap iterations.
At each iteration, a randomly selected boundary tract is moved from its
current district to an adjacent district, if:
(a) the move reduces the maximum deviation from ideal population, and
(b) both districts remain non-empty and contiguous after the swap.
The phase terminates early if all districts satisfy the balance constraint.
\end{definition}

The rebalance phase is identical to the boundary-swap logic used in the
MultiScale structure-layer algorithm, ensuring consistency across the
platform.

\subsection{Contiguity Guarantee}

\begin{proposition}[Contiguity by Construction]
\label{prop:contiguity}
Let $G_v$ be connected.
The BFS growth phase (Definition~\ref{def:growth}) produces an assignment
$A$ in which every district $d$ induces a connected subgraph.
\end{proposition}

\begin{proof}
By induction on the number of tracts assigned to district $d$.

\textbf{Base case.} After initialisation, district $d$ contains exactly
$\{\mathrm{seed}_d\}$, which is trivially connected.

\textbf{Inductive step.} Suppose district $d$ currently induces a connected
subgraph $S_d$.
The next tract $i$ assigned to $d$ is popped from $\mathcal{Q}$ with
assignment label $d$; by Definition~\ref{def:growth} step~3(d),
$i$ was pushed to $\mathcal{Q}$ because it is a neighbour of some tract
$j \in S_d$ (specifically the tract that was assigned to $d$ when $i$ was
first encountered).
Therefore $i$ is adjacent to $j \in S_d$, and $S_d \cup \{i\}$ is
connected.

By induction, $S_d$ is connected at every stage of growth.
After all tracts are assigned, each district is a connected subgraph.
\end{proof}

\begin{remark}
The skip at step~3(b) (``already assigned'') is critical: it is possible
for the same tract $i$ to appear in $\mathcal{Q}$ multiple times, each
with a different district label, if $i$ is adjacent to tracts from
multiple districts.
The first pop assigns $i$ to the district whose frontier reached $i$
first (i.e., the district with the lowest current population at the time
of the push).
Subsequent pops for $i$ are discarded.
This does not break the contiguity argument because the proving induction
only requires that the assigned tract is adjacent to some previously
assigned tract in the same district---a condition guaranteed by the push
at step~3(d).
\end{remark}

\subsection{Computational Complexity}

\begin{proposition}[Runtime Complexity]
\label{prop:complexity}
Let $n = |V_v|$ and $m = |E_v|$.
The \bfsg{} algorithm (seed selection + BFS growth + rebalance) runs in
$O((n + m) \log n)$ time.
For redistricting graphs, which are planar with $m = O(n)$, this simplifies
to $O(n \log n)$.
\end{proposition}

\begin{proof}
\textit{Seed selection.}
Each BFS-farthest computation requires one BFS from all current seeds
simultaneously, which costs $O(n + m)$.
With $k$ seeds, the total seed selection cost is $O(k(n + m))$.
For redistricting bisection $k = 2$; for larger $k$, $k \ll n$.

\textit{BFS growth.}
Each of the $n$ tracts is pushed to $\mathcal{Q}$ at most once per
adjacent district (at most $\deg(i)$ times).
Each push/pop is $O(\log n)$.
Total: $O((n + m) \log n)$.

\textit{Rebalance.}
The rebalance phase performs at most 200 boundary-swap checks, each
requiring a contiguity BFS $O(n + m)$.
Total: $O(200(n + m)) = O(n + m)$.

\textit{Overall.}
$O((n + m) \log n) + O(n + m) = O((n + m) \log n)$.
For planar $m = O(n)$: $O(n \log n)$.
\end{proof}

\begin{remark}
The $O(n \log n)$ bound matches a single \metis{} call.
Empirically, \bfsg{} is slightly faster than \metis{} in practice because
\metis{}'s constant factors include memory allocation for the multilevel
coarsening data structures, whereas \bfsg{}'s inner loop is a simple
priority-queue operation with a single population accumulator per district.
\end{remark}

\subsection{Formal Pseudocode}

Algorithm~\ref{alg:bfs} gives the complete \bfsg{} bisection procedure.

\begin{algorithm}[t]
\caption{\bfsg{}-\textsc{Bisect}($G_v$, $\mathtt{base\_seed}$, $\mathtt{balance\_post\_iters}$)}
\label{alg:bfs}
\begin{algorithmic}[1]
\State $n \gets |V_v|$; \quad $P \gets \sum_{i \in V_v} p_i$ \Comment{total subgraph population}
\State $s \gets \mathrm{BFS\text{-}Seed}(\mathtt{base\_seed})$ \Comment{Definition~\ref{def:bfs-seed}}
\State $\mathrm{rng} \gets \mathrm{SmallRng::seed\_from\_u64}(s)$
\State $\mathrm{seed}_0 \gets \mathrm{WeightedSample}(V_v,\; \{p_i\},\; \mathrm{rng})$
       \Comment{population-weighted}
\State $D_0 \gets \mathrm{BFS\text{-}Distances}(\mathrm{seed}_0,\; G_v)$
\State $\mathrm{seed}_1 \gets \arg\max_{i \in V_v} D_0[i]$
       \Comment{k-farthest}
\State $A \gets [\bot;\; n]$; \quad
       $A[\mathrm{seed}_0] \gets 0$; \quad $A[\mathrm{seed}_1] \gets 1$
\State $\mathrm{cur}[0] \gets p_{\mathrm{seed}_0}$; \quad
       $\mathrm{cur}[1] \gets p_{\mathrm{seed}_1}$
       \Comment{current district populations}
\State $\mathcal{Q} \gets \emptyset$ \Comment{min-heap keyed on cur[district]}
\For{$d \in \{0, 1\}$}
  \For{$j \in \mathrm{Neighbours}(\mathrm{seed}_d,\; G_v)$}
    \If{$A[j] = \bot$}
      \State $\mathcal{Q}.\mathrm{push}(\mathrm{cur}[d],\; j,\; d)$
    \EndIf
  \EndFor
\EndFor
\While{$\mathcal{Q} \neq \emptyset$}
  \State $(\text{---},\; i,\; d) \gets \mathcal{Q}.\mathrm{pop\_min}()$
  \If{$A[i] \neq \bot$} \textbf{continue} \EndIf \Comment{already assigned}
  \State $A[i] \gets d$; \quad $\mathrm{cur}[d] \mathrel{+}= p_i$
  \For{$j \in \mathrm{Neighbours}(i,\; G_v)$}
    \If{$A[j] = \bot$}
      \State $\mathcal{Q}.\mathrm{push}(\mathrm{cur}[d],\; j,\; d)$
    \EndIf
  \EndFor
\EndWhile
\State $A \gets \mathrm{Rebalance}(A,\; G_v,\; \{p_i\},\; 0.005,\; \mathtt{balance\_post\_iters})$
\State \textbf{return} $(\{i : A[i] = 0\},\; \{i : A[i] = 1\})$
\end{algorithmic}
\end{algorithm}

\subsection{CLI Invocation and Implementation}

\bfsg{} bisection is exposed via the \texttt{--structure bfs-growth}
flag in the \texttt{redist} binary.
The full invocation for North Carolina is:

\begin{verbatim}
redist state --state NC --year 2020 --structure bfs-growth
\end{verbatim}

No additional structure-layer parameters are required.
The \texttt{--balance-tolerance} flag (compositor-level) controls the
rebalance phase tolerance (default $0.5\%$).
The YAML configuration is:

\begin{verbatim}
algorithm:
  structure: bfs-growth
  weights: geographic
  search: single
  balance_tolerance: 0.5
workers: 8
years: ["2020"]
\end{verbatim}

When combined with \texttt{--search convergence --convergence-threshold 600},
the SeedCompositor runs \cs{} over BFS growth bisections, providing
multi-seed quality improvement.
